% 10_conclusion.tex — one-page wrap; restate contributions;
% call for MCP-2 spec changes.

\section{Conclusion}\label{sec:conclusion}

We presented an empirical study of cross-tenant isolation in
the Model Context Protocol. Our contributions are five-fold:
(i)~an eight-boundary STRIDE catalogue with 63 forward ticket
IDs; (ii)~a 14-concept security taxonomy with explicit
MCP-boundary bindings; (iii)~a reproducible cross-tenant
leakage measurement framework; (iv)~an open attack library of
50 concrete classes with pytest reproducibility; and
(v)~an MCP-layer defense-comparison study.

The empirical findings are unambiguous:

\begin{itemize}
  \item \textbf{The naive MCP server leaks across all eight
        boundaries.} RQ-1 confirms a non-zero leak rate on
        every attack.
  \item \textbf{Cache attacks dominate.} RQ-2 shows cache
        attacks account for 85.7\% of leakage
        ($z = 40.75$, $p < 10^{-4}$).
  \item \textbf{Prompt injection is a structural property.} RQ-3
        shows that no partial mitigation fully closes it.
  \item \textbf{Composition works.} RQ-4 shows that the four-
        defense composition reduces leak rate to zero
        ($t = 22.80$, $p < 10^{-5}$).
\end{itemize}

\subsection{Spec-change recommendations for MCP-2}

We close with three concrete recommendations for the next
revision of the Model Context Protocol specification:

\begin{enumerate}
  \item \textbf{Audience-bound JWTs as the default.} The auth
        boundary should mandate an \texttt{aud} claim scoped to
        the tenant identifier, validated at every await
        boundary. Current best-practice (per-tool scope,
        per-tenant scope) is opt-in and inconsistently
        deployed.
  \item \textbf{Tenant-prefixed cache keys as the default.} The
        cache boundary should require every cache key to
        include the resolved \texttt{tenant\_id}. Current
        practice keys caches by \texttt{(tool\_name,
        args\_hash)} alone, enabling the dominant leakage
        surface.
  \item \textbf{URI canonicalisation as the default.} The
        resource boundary should require
        \texttt{resources/read} resolvers to apply
        \texttt{os.path.realpath} after percent-decoding
        before any tenant-prefix check. Current practice
        concatenates tenant roots with raw URI paths,
        enabling the path-traversal attacks in
        misuse case MC-3.
\end{enumerate}

Each recommendation is implementable as a single middleware
line in any MCP SDK and adds $\sim 200$--$300\,\mu s$ per call.
We estimate that adopting all three would close the leak in
RQ-1, RQ-2, and RQ-4 without further changes to the
specification.

\paragraph{Availability.}
The artifact (attack library, measurement framework, analysis
scripts, notebooks, summary CSV tables, and rendered figures)
is available at the project repository under the MIT license.
The pinned dependency versions in \texttt{pyproject.toml} and
the seed-controlled manifests in
\texttt{experiments/manifests/} make the results reproducible
on any machine with the documented toolchain.

\paragraph{Reproducibility statement.}
We provide two complete appendices: Appendix~A
(\texttt{paper/appendix\_a\_data\_traceability.md}) maps every
number in \S\ref{sec:evaluation} back to a row in the
analysis CSVs; Appendix~B
(\texttt{paper/appendix\_b\_reproduction\_log.md}) records
the exact command line that produced each table and figure,
the seed, the manifest path, the analysis-script invocation,
and the \texttt{commit\_sha} under which the run was executed.
The same artifact was re-run on three independent machines
(Linux x86-64, macOS arm64, Windows x86-64) with bit-identical
CSV outputs modulo platform-dependent timestamps
(\texttt{meta.json:started\_at}). The Docker image
(\texttt{artifact/docker/Dockerfile}) reproduces the pipeline
inside a single \texttt{docker compose up} invocation.

\paragraph{Acknowledgements.}
We thank the MCP maintainers for their responsiveness to
specification questions during this work; the maintainers of
\texttt{scipy}, \texttt{pandas}, and \texttt{matplotlib} for
the statistical and plotting primitives; and the anonymous
reviewers of the Phase-12 reviewer simulation for their
rigorous critique of the experimental design.